\documentclass[12pt,a4paper,oneside]{report}
\usepackage[margin=25mm,headheight=20pt]{geometry}
\usepackage{fontspec}
\setmainfont{THSarabunPSK}[Path=assets/,Extension=.ttf,UprightFont=*-Regular,BoldFont=*-Bold,ItalicFont=*-Italic,BoldItalicFont=*-BoldItalic]
\setsansfont{THSarabunPSK}[Path=assets/,Extension=.ttf,UprightFont=*-Regular,BoldFont=*-Bold]
\setmonofont{lmmono10-regular.otf}[Scale=MatchLowercase]
% TH SarabunPSK: body 16pt, titles 18pt, small (header, diagram labels) 14pt
% sizes are Word/PDF points (bp): 16bp = 16.06 TeX pt
\makeatletter
\renewcommand\normalsize{\@setfontsize\normalsize{16.06}{20}%
  \abovedisplayskip 10\p@ \@plus3\p@ \@minus6\p@ \belowdisplayskip \abovedisplayskip
  \abovedisplayshortskip \z@ \@plus3\p@ \belowdisplayshortskip 6\p@ \@plus3\p@ \@minus3\p@}
\renewcommand\small{\@setfontsize\small{14.05}{17}}
\makeatother
\newcommand{\titlesize}{\fontsize{18.07}{22}\selectfont}
\normalsize
\XeTeXlinebreaklocale "th"
\XeTeXlinebreakskip=0pt plus 1pt
\usepackage{pifont,capt-of,graphicx,booktabs,tabularx,longtable,array,amsmath,enumitem,xcolor,tikz,fancyhdr,titlesec,float}
\usetikzlibrary{arrows.meta,positioning,shapes.geometric}
\usepackage[unicode,colorlinks=true,linkcolor=black,urlcolor=black,citecolor=black]{hyperref}
\usepackage{bookmark}
\definecolor{navy}{HTML}{000000}
\definecolor{pale}{HTML}{F2F2F2}
\renewcommand{\contentsname}{สารบัญ}
\renewcommand{\listtablename}{สารบัญตาราง}
\renewcommand{\listfigurename}{สารบัญภาพ}
\renewcommand{\tablename}{ตารางที่}
\renewcommand{\figurename}{ภาพที่}
\renewcommand{\bibname}{บรรณานุกรม}
\renewcommand{\chaptername}{บทที่}
\renewcommand{\appendixname}{ภาคผนวก}
\titleformat{\chapter}[display]{\bfseries\color{navy}\titlesize}{\chaptername\ \thechapter}{5pt}{\titlesize}
\titlespacing*{\chapter}{0pt}{0pt}{22pt}
\titleformat{\section}{\normalsize\bfseries\color{navy}}{\thesection}{0.6em}{}
\titleformat{\subsection}{\normalsize\bfseries}{\thesubsection}{0.6em}{}
\setlength{\parindent}{1.3em}
\setlength{\parskip}{5pt}
\linespread{1}
\setlength{\emergencystretch}{3em}
\tolerance=3000
\setlist{itemsep=3pt,topsep=5pt,leftmargin=1.6em}
\renewcommand{\arraystretch}{1.35}
\newcolumntype{Y}{>{\raggedright\arraybackslash}X}
\pagestyle{fancy}
\fancyhf{}
\fancyhead[L]{\small HENJAI | เห็นใจ}
\fancyhead[R]{\small ทีม Rushhour}
\fancyfoot[C]{\normalsize\thepage}
\renewcommand{\headrulewidth}{0.3pt}
\fancypagestyle{plain}{\fancyhf{}\fancyfoot[C]{\normalsize\thepage}\renewcommand{\headrulewidth}{0pt}}
\DeclareRobustCommand{\file}[1]{\nolinkurl{#1}}
\newcommand{\note}[1]{\par\noindent\colorbox{pale}{\parbox{\dimexpr\linewidth-2\fboxsep}{#1}}\par}
\newcommand{\frontchapter}[1]{\chapter*{#1}\addcontentsline{toc}{chapter}{#1}}
\newcommand{\projectcover}[1]{%
\hypersetup{pageanchor=false}\begin{titlepage}\centering\titlesize\bfseries
HENJAI\par
\vspace{4mm}เห็นใจ\par
\vspace{9mm}#1\par
\vspace{7mm}โครงการผู้ช่วยบรรยายภาพด้วยปัญญาประดิษฐ์\\ที่เข้าใจอารมณ์และบริบทสำหรับผู้พิการทางสายตา\par
\vspace{3mm}Emotion \& Context-Aware AI Vision Assistant\par
\vfill
ผู้จัดทำ\par\vspace{4mm}
นายกิตติพงศ์ สมสอาด\\
นายชินพัฒน์ ขุนเจริญ\\
นายเรข์ณพัศ สุวรรณปริสุทธิ์\par
\vspace{4mm}ทีม Rushhour\\นิสิตชั้นปีที่ 2 หลักสูตรวิศวกรรมคอมพิวเตอร์และเทคโนโลยีดิจิทัล (CEDT)\\ภาควิชาวิศวกรรมคอมพิวเตอร์\\คณะวิศวกรรมศาสตร์ จุฬาลงกรณ์มหาวิทยาลัย\par
\end{titlepage}\hypersetup{pageanchor=true}}

\hypersetup{pdftitle={HENJAI คู่มือการใช้งานระบบ},pdfauthor={ทีม Rushhour}}
\begin{document}
\projectcover{คู่มือการใช้งานระบบ\\User Manual}
\pagenumbering{roman}
\frontchapter{เกี่ยวกับคู่มือนี้}
คู่มือนี้อธิบายวิธีเตรียม เริ่มใช้งาน และแก้ปัญหาเบื้องต้นของ HENJAI สำหรับผู้ใช้ ผู้ดูแล และผู้สาธิต ครอบคลุมต้นแบบเว็บกับต้นแบบบนเครื่องคอมพิวเตอร์ ขั้นตอนติดตั้งสำหรับนักพัฒนาและรายละเอียด API อยู่ในคู่มือการพัฒนาระบบแยกเล่ม

ระบบเป็นต้นแบบสำหรับบรรยายภาพผ่านเสียงภาษาไทย โดยวิเคราะห์สิ่งที่กล้องเห็นในขณะนั้น ไม่ได้ยืนยันอารมณ์ภายในบุคคล ไม่ได้ทราบสิ่งที่อยู่นอกมุมกล้อง และยังไม่มีผลรับรองการนำทางอย่างอิสระ การสาธิตควรทำในบริเวณที่ควบคุมได้ โดยมีผู้ดูแลช่วยตรวจสิ่งแวดล้อม

เอกสารจัดทำจากไฟล์โครงการ ณ วันที่ 20 ตุลาคม 2569 ยังไม่ได้ทดสอบกล้อง เสียง หรือบริการออนไลน์ใหม่ในรอบจัดทำเอกสาร ชื่อรุ่นโมเดลในซอร์สเป็นค่าที่มีอยู่ในโครงการ ผู้ดูแลต้องตรวจความพร้อมของบัญชีบริการก่อนวันใช้งาน

\tableofcontents
\clearpage\pagenumbering{arabic}
\chapter{รู้จักระบบและเตรียมใช้งาน}
\section{เลือกเวอร์ชันที่เหมาะกับการสาธิต}
\begin{table}[H]\centering\caption{ความแตกต่างของระบบสองรูปแบบ}
\begin{tabularx}{\linewidth}{lYY}\toprule ประเด็น & เว็บ & บนเครื่องคอมพิวเตอร์\\\midrule
ผู้เตรียมระบบ & ผู้ดูแลเผยแพร่เว็บและตั้งค่า key & ผู้ดูแลติดตั้ง Python และโมเดล\\
กล้อง & กล้องของมือถือหรือเครื่องที่เปิดเว็บ & กล้องที่เชื่อมกับเครื่องเซิร์ฟเวอร์\\
เสียง & เสียงภาษาไทยของเบราว์เซอร์ & gTTS หรือเสียง Kanya บน macOS\\
อินเทอร์เน็ต & ต้องใช้สำหรับคำบรรยาย Gemini & ต้องใช้สำหรับคำบรรยาย Gemini\\
เตือนระยะในเครื่อง & ไม่มีชุดโมเดลสายเร็ว & มีเมื่อเปิด fast และโหลดโมเดลครบ\\
หน้าจอตรวจสอบ & เพิ่ม \file{?dev=1} หลัง URL & เปิด \file{/dev}\\\bottomrule
\end{tabularx}\end{table}

คู่มือนี้ไม่ระบุ URL สาธารณะที่ยังไม่ได้ตรวจว่าเปิดใช้งานอยู่ ให้ผู้ดูแลแจ้ง URL ของ deployment ที่เตรียมไว้จริง การเปิดไฟล์ HTML จากโฟลเดอร์โดยตรงไม่สามารถแทนเซิร์ฟเวอร์ API ได้

\section{สิ่งที่ต้องเตรียม}
\begin{enumerate}
\item อุปกรณ์มีกล้องที่มองเห็นภาพชัดเจน และลำโพงหรืออุปกรณ์ฟังเสียงที่เหมาะกับสถานการณ์
\item อินเทอร์เน็ตสำหรับการวิเคราะห์บนคลาวด์ และบัญชีบริการที่ผู้ดูแลตรวจ key กับโควตาแล้ว
\item เว็บเบราว์เซอร์ที่อนุญาตใช้กล้อง มีเสียงภาษาไทย และสามารถเปิด URL ที่ปลอดภัยได้
\item พื้นที่สาธิตที่ไม่มีความเสี่ยงจากรถ บันได หรือสิ่งกีดขวางที่ผู้ใช้ต้องหลบด้วยตัวเอง
\end{enumerate}

\section{ตรวจความพร้อมก่อนเริ่ม}
ให้ผู้ดูแลตรวจกล้องด้วยฉากง่าย ๆ เช่น โต๊ะและเก้าอี้ในห้อง ตรวจว่าทิศซ้ายขวาตรงกับมุมมองผู้ใช้และไม่ได้กลับภาพโดยไม่ตั้งใจ จากนั้นทดสอบเสียงหนึ่งประโยค เปิดระดับเสียงพอได้ยินแต่ยังรับรู้เสียงรอบตัวได้ และแจ้งให้ผู้ที่อยู่ในบริเวณทราบว่าจะมีการส่งภาพไปประมวลผลภายนอก

หลีกเลี่ยงภาพเอกสารลับ ข้อมูลส่วนตัว และหน้าจอที่มีรหัสผ่าน ในโหมดบรรยายคลาวด์ ภาพที่เลือกกับบริบทบางส่วนจะถูกส่งไปยังบริการ AI การกดหยุดหยุดการรับภาพและคำขอรอบถัดไป แต่ไม่ย้อนกลับคำขอที่ส่งไปแล้ว

\chapter{การใช้งานผ่านเว็บ}
\section{เปิดระบบและอนุญาตอุปกรณ์}
\begin{enumerate}
\item เปิด URL ที่ผู้ดูแลเตรียมไว้ผ่าน HTTPS บนอุปกรณ์ที่ต้องการใช้กล้อง
\item ตรวจว่าระดับเสียงเปิดอยู่และไม่ได้ส่งเสียงไปยังอุปกรณ์ Bluetooth ที่อยู่อีกที่หนึ่ง
\item กดปุ่มหลักเพื่อเริ่ม ระบบจะขอสิทธิ์กล้อง ให้เลือกอนุญาต หากไม่ปรากฏหน้าต่าง ให้ตรวจสิทธิ์ของเว็บไซต์ในเบราว์เซอร์
\item ถือหรือติดกล้องให้ทิศภาพสัมพันธ์กับสิ่งที่ผู้ใช้ต้องการรับรู้ เว็บร้องขอกล้องด้านหลังเป็นค่าเริ่มที่ต้องการ แต่เครื่องอาจเลือกอุปกรณ์อื่นตามที่มี
\item รอเสียงบรรยายแรก การเริ่มต้นอาจใช้เวลาจากเครือข่ายและบริการ AI
\end{enumerate}

\section{ฟังคำบรรยายและเปลี่ยนมุมกล้อง}
ระบบพยายามบรรยายสิ่งสำคัญอย่างกระชับและพูดเมื่อมีข้อมูลใหม่ เมื่อภาพนิ่งอาจไม่มีเสียงพูดเพิ่มเติม ซึ่งเป็นพฤติกรรมที่ออกแบบไว้ หากต้องการเปลี่ยนสิ่งที่รับรู้ ให้หันกล้องอย่างช้า ๆ ไปยังบริเวณใหม่และรอให้ระบบวิเคราะห์รอบถัดไป ไม่จำเป็นต้องกดเริ่มซ้ำทุกภาพ

ถ้าคำบรรยายไม่ตรงกับสิ่งที่ผู้ดูแลเห็น ให้หยุดและปรับมุมกล้องหรือแสง อย่าอาศัยข้อความนั้นตัดสินใจเคลื่อนไปยังพื้นที่เสี่ยง ข้อความเกี่ยวกับสีหน้าเป็นการตีความจากภาพเดียว จึงอาจผิดในกรณีหาว จาม หน้าหัน หรือแสงเงา

\section{การอ่านเอกสาร}
นำเอกสาร ป้าย หรือข้อความให้กินพื้นที่ส่วนใหญ่ของภาพ ตั้งกล้องให้นิ่ง จัดแสงให้ตัวอักษรชัดและไม่มีแสงสะท้อน เมื่อระบบตรวจว่าเป็นเอกสาร จะอ่านข้อความที่ได้หลังคำเกริ่น หากเอกสารยาว ให้แบ่งทีละส่วนและตรวจความครบถ้วนกับผู้ดูแล เพราะคำตอบจากโมเดลมีขนาดจำกัดและยังไม่มีผลทดสอบการอ่านครบทุกประเภทเอกสาร

หากฟังข้อความที่มีผลต่อเงิน ยา หรือข้อมูลสำคัญ ให้ตรวจจากต้นฉบับด้วยผู้ช่วยหรือเครื่องมือที่เหมาะสมก่อนดำเนินการ ไม่สรุปจากการที่ระบบอ่านออกเสียงได้ว่าตัวอักษรถูกต้องครบทุกตัว

\section{หยุดและจบการใช้งาน}
กดปุ่มหลักอีกครั้งเพื่อหยุด เว็บจะหยุด track กล้องและยกเลิกเสียงที่กำลังพูด ปิดหน้าเว็บเมื่อเสร็จ หากต้องการเริ่มบริบทใหม่ ให้รีโหลดหน้าแล้วเริ่มอีกครั้ง ความจำในเว็บอยู่ในหน้านั้น ไม่ได้ออกแบบเป็นประวัติระยะยาวของผู้ใช้

\section{หน้าจอตรวจสอบสำหรับผู้ดูแล}
เติม \file{?dev=1} ต่อท้าย URL เพื่อดูข้อมูลโหมด เวลา และความเปลี่ยนแปลงของภาพ ใช้ตรวจการสาธิตและช่วยอธิบายการข้ามภาพนิ่ง ข้อมูลเวลาที่แสดงมาจากการเรียกวิเคราะห์ ไม่ใช่เวลารวมจนผู้ใช้ได้ยินคำบรรยายจบ

\chapter{การใช้งานบนเครื่องคอมพิวเตอร์}
\section{เปิดบริการโดยผู้ดูแล}
หลังติดตั้งตามคู่มือพัฒนา ให้เปิด terminal ใน \file{source_files/poc/} แล้วรันคำสั่ง
\begin{verbatim}
python server.py
\end{verbatim}
เปิด \url{http://127.0.0.1:8000/} สำหรับหน้าจอผู้ใช้ หรือ \url{http://127.0.0.1:8000/dev} สำหรับผู้ดูแล บริการนี้ใช้กล้องของเครื่องที่รัน Python ไม่ใช่กล้องของอุปกรณ์ที่เปิดหน้าเว็บจากที่อื่น

\section{ตั้งกล้องและเริ่มต้น}
\begin{enumerate}
\item เปิด \file{/dev} และเลือกกล้องที่มีภาพจริง ตรวจภาพตัวอย่างก่อนกดเริ่ม
\item ตรวจค่า API key ที่ตั้งไว้ในเครื่อง และเลือกโมเดลที่บัญชีของผู้ดูแลใช้งานได้
\item หากสลับกล้องแล้วภาพไม่ขึ้น ให้หยุดและเริ่มใหม่ พร้อมตรวจว่าสิทธิ์กล้องของโปรแกรม Python ได้รับอนุญาต
\item ตั้งกล้องให้นิ่งในช่วงเริ่มต้นเพื่อให้ระบบประมาณ noise สำหรับเกณฑ์ภาพเปลี่ยน
\item ไปหน้าจอผู้ใช้หลักแล้วกดเริ่ม ฟังคำบรรยายและสถานะจากระบบ
\end{enumerate}

\section{ค่าที่ผู้ดูแลอาจปรับ}
\begin{table}[H]\centering\caption{ค่าควบคุมต้นแบบบนเครื่อง}
\begin{tabularx}{\linewidth}{lY}\toprule รายการ & แนวทางใช้\\\midrule
กล้อง & เลือกอุปกรณ์ที่ให้ภาพตามมุมมองผู้ใช้\\
Mirror & ตรวจซ้ายขวากับวัตถุจริงก่อนเปลี่ยนค่า\\
Every & ช่วงรอระหว่างรอบ ไม่ใช่คำรับรองเวลาตอบสนอง\\
Still / Scene / Alert & เกณฑ์ภาพเปลี่ยน; ปกติใช้ auto-calibration\\
Recalibrate & สอบเทียบ noise ใหม่เมื่อเปลี่ยนกล้องหรือการติดตั้ง\\
Forget & ล้างประวัติคำบรรยายและสรุปใน session\\
Fast & เปิดเส้นทางโมเดลในเครื่อง ต้องเตรียม dependency และโมเดลครบ\\\bottomrule
\end{tabularx}\end{table}

การตั้งเกณฑ์เองจะปิดโหมดอัตโนมัติในซอร์ส หากไม่มั่นใจให้สอบเทียบใหม่แทนการปรับทุกค่า การเพิ่มความไวมากเกินไปอาจเพิ่มคำขอคลาวด์และทำให้ระบบพูดถี่ขึ้น

\section{การเตือนระยะใกล้}
เส้นทางในเครื่องอาจเตือนเมื่อค่าระยะทำนายต่ำกว่า 0.5 เมตรและยังไม่เคยเตือนในรอบนั้น แล้วรอจนระยะมากกว่า 0.9 เมตรจึงพร้อมเตือนใหม่ การไม่มีเสียงไม่ได้หมายความว่าทางข้างหน้าปลอดภัย ระบบไม่ได้เตือนทุกวัตถุหรือทุกการเคลื่อนไหว และอาจระงับการพูดเมื่อมีเสียงอื่นกำลังเล่น

\section{หยุดระบบ}
กดหยุดในหน้าจอก่อน แล้วใช้ Ctrl+C ใน terminal หากจะปิดบริการ หากใช้ iPhone เป็น Continuity Camera ให้ตรวจว่าไม่ได้ล็อกจอหรือหลุดการเชื่อมต่อระหว่างการสาธิต

\chapter{แก้ปัญหาและตรวจรับเบื้องต้น}
\section{อาการที่พบบ่อย}
\begin{longtable}{p{36mm}p{110mm}}
\caption{แนวทางแก้ปัญหาเบื้องต้น}\\\toprule อาการ & สิ่งที่ตรวจและการดำเนินการ\\\midrule\endfirsthead
\toprule อาการ & สิ่งที่ตรวจและการดำเนินการ\\\midrule\endhead
ไม่เห็นภาพกล้อง & ตรวจสิทธิ์เว็บไซต์หรือ Python ปิดแอปที่ใช้กล้องอยู่ ตรวจว่าเปิด HTTPS บนเว็บ และเลือกกล้องให้ถูก\\
มีภาพแต่ไม่มีเสียง & ตรวจระดับเสียง อุปกรณ์ส่งเสียง และเสียงภาษาไทยในระบบ; กดเริ่มด้วยผู้ใช้เพื่อปลดล็อกเสียงเบราว์เซอร์\\
ระบบเงียบเมื่อกล้องนิ่ง & อาจเป็นการข้ามภาพนิ่งหรือกันซ้ำ ให้เปลี่ยนฉากช้า ๆ และดูสถานะจากหน้าตรวจสอบ\\
ขึ้น 429 / quota & หยุดสาธิตและให้ผู้ดูแลตรวจโควตา ไม่กดเริ่มซ้ำถี่ ๆ\\
ขึ้น 404 / model & ให้ผู้ดูแลตรวจรุ่นโมเดลที่เปิดใช้งานสำหรับบัญชี แล้วตั้งค่าโมเดลใหม่ตามคู่มือพัฒนา\\
หมดเวลา / ต่อเน็ตไม่ได้ & ตรวจเครือข่าย รอให้เสถียร แล้วทดลองใหม่กับฉากง่าย ๆ\\
สายเร็วขึ้น error & ตรวจ dependency และไฟล์โมเดล ให้ผู้ดูแลแก้ก่อนเปิด fast อีกครั้ง\\
ภาพดำจาก iPhone & ปลดล็อก iPhone รอสัญญาณ หรือเลือกกล้องอื่น แล้วหยุดและเริ่มใหม่\\
คำบรรยายไม่ตรงภาพ & ตรวจมุมกล้อง แสง และสิ่งบัง ล้างความจำถ้าบริบทเดิมค้าง บันทึกกรณีให้ผู้พัฒนาตรวจ\\\bottomrule
\end{longtable}

\section{รายการตรวจรับสำหรับการสาธิต}
\begin{enumerate}
\item เริ่มและหยุดระบบได้ กล้องหยุดเมื่อกดหยุด และหน้าเว็บไม่มีสถานะผิดพลาดค้าง
\item ได้ยินเสียงภาษาไทยและปรับระดับเสียงได้
\item ฉากง่าย ๆ ให้คำบรรยายที่ผู้ดูแลตรวจเทียบได้
\item การเปลี่ยนฉากทำให้เกิดข้อมูลใหม่ และภาพนิ่งไม่ทำให้พูดวนโดยไม่จำเป็น
\item ทดสอบอ่านข้อความสั้นแล้วตรวจเทียบตัวอักษร ไม่ใช้ผลนี้รับรองเอกสารยาว
\item จำลองเครือข่ายหรือสิทธิ์กล้องไม่พร้อมและตรวจว่าผู้ดูแลเข้าใจวิธีกู้คืน
\end{enumerate}
รายการนี้เป็นแนวทางตรวจรับที่เสนอ ยังไม่ใช่บันทึกว่าระบบผ่านทุกข้อในเครื่องปัจจุบัน ให้ผู้ดูแลบันทึกวันที่ อุปกรณ์ เบราว์เซอร์ โมเดล และผลจริงทุกครั้ง

\section{แจ้งกรณีให้ผู้พัฒนา}
บันทึกขั้นตอนที่ทำ ข้อความผิดพลาด รุ่นอุปกรณ์ เวลาที่เกิด และรูปแบบการใช้งาน หากจำเป็นต้องส่งภาพให้ขออนุญาตและตัดข้อมูลส่วนตัวก่อน ห้ามแนบ API key หรือไฟล์ \file{.env} ในข้อความแจ้งปัญหา
\end{document}
